\scanpage{60}\mspage{59}
J'en viens à \unsure{des|du} \uline{\unsure{axiomes|axiome} de disjonction}
(At 8 était l' ``axiome des lieux'').

Notons d'abord

\deleted{\uline{At} 9} \deleted{Soient $\unsure{X|F}$}

\uline{Proposition} Soit $\unsure{\mathcal F|F}$ un atelier
At 1 -- At 3 (on \unsure{ne|n'y} va pas utiliser $\ll$).
Soient $F$ et $G\in\mathcal F$, alors $\unsure{\overline{1}|T|\mathrm{I}}$
% slanted
\margin{\unsure{Remarquons|Remarquer} \unsure{tt de suite} que $X\parallel Y\Longrightarrow\operatorname{Omb}(X)\cap\operatorname{Omb}Y=\varnothing$ a fortiori $\operatorname{omb}(X)\cap\operatorname{omb}(Y)=0$}
\[
F\parallel G\quad\Longleftrightarrow\quad\forall X\in\widetilde F,Y\in\widetilde G,\qquad X\parallel Y.
\]
Immédiat. \unsure{Egalmt|Egalement}, \deleted{\unsure{on trouve}} \inserted{on voudrait}

\begingroup\makeatletter\tagsleft@true\makeatother
\[
\begin{aligned}
&\text{\deleted{\uline{Proposition}} \deleted{\unsure{Corollaire}} Si }\ F\parallel G,\quad F'\ll F,\ G'\ll G,\\
&\text{alors }\ F'\parallel G'.
\end{aligned}
\tag*{(1.58)}
\]
\endgroup

En effet, \deleted{\unsure{il suffit|il soit}} on est ramené au cas où
$F'=X',\ G'=Y'$ \inserted{(\unsure{par prop})}, puis au cas
$F=X,\ G=Y$, en utilisant At 4.
% slanted
\margin{\uline{NB} Le seul pb, c'est de prouver que $F'\mathrel{\tikz[baseline=-.45ex,x=1em,y=1em]{\draw[line width=.055em] (.52,.27)--(.05,.055)--(.48,-.11); \draw[line width=.055em] (.14,.21)--(.46,.04)--(0,-.17);}}G'$ \deleted{\unsure{quand}} \deleted{\unsure{une fois cela vu}}, \inserted{\unsure{c'est ``oui''|c'est ``que''}}, $F'\overset{-\,\searrow}{\parallel}G'$. \deleted{\unsure{c'est ici|c'est lui}} Car sinon \unsure{car, un|car, le} $Z\in M$ \deleted{$Z$} avec $Z\leq F',\ Z\leq G'$ donne $Z\ll F,\ Z\ll G$ Or \deleted{dans} $\unsure{\widetilde{F\vee G}|\widehat{F\vee G}}$ les \unsure{4|2} relations $\ll$ et $\leq$ \unsure{coïncident}, de $\parallel$ \unsure{c.a.d.} le $\operatorname{Inf}$ pour $\leq$ $\unsure{\nexists|\neq}$, et pour $\ll$ on trouve absurdité\ldots}

Même, dans le cas ensembliste à figures de t.f.
(et $=$ finies) cela ne semble pas résulter
des autres axiomes -- cf page 55,
\unsure{Atens tf|Atens t.f|Atensfig} 3 nous dit quelles conditions
mettre à la relation $\underset{M}{\mathrel{\tikz[baseline=-.45ex,x=1em,y=1em]{\draw[line width=.055em] (.52,.27)--(.05,.055)--(.48,-.11); \draw[line width=.055em] (.14,.21)--(.46,.04)--(0,-.17);}}}$ \deleted{\unsure{de pré}} : d'être
stable par passage aux sous-strates, et
d' \deleted{être} plus forte que la relation $\underset{\mathrm{ens}}{\mathrel{\tikz[baseline=-.45ex,x=1em,y=1em]{\draw[line width=.055em] (.52,.27)--(.05,.055)--(.48,-.11); \draw[line width=.055em] (.14,.21)--(.46,.04)--(0,-.17);}}}$.
